% !TeX program = xelatex
% 图1：发布审查方法总体流程（三阶段，分别对应 3.2 / 3.3 / 3.4）
% 阶段名置于虚线框内的标题栏（不带节号）；子步骤仅保留主名称；第三阶段为一个总阶段
% 独立编译，用于导出高分辨率位图
\documentclass[border=2pt,varwidth=false]{standalone}

\usepackage{fontspec}
\usepackage{xeCJK}
\usepackage{amsmath,amssymb,bm}
\usepackage{tikz}
\usetikzlibrary{arrows.meta,positioning,fit,backgrounds,calc}

\IfFontExistsTF{Songti SC}
  {\setCJKmainfont{Songti SC}}
  {\IfFontExistsTF{SimSun}{\setCJKmainfont{SimSun}}
    {\setCJKmainfont{FandolSong-Regular.otf}[BoldFont=FandolSong-Bold.otf]}}
\IfFontExistsTF{Heiti SC}
  {\setCJKsansfont{Heiti SC}}
  {\IfFontExistsTF{SimHei}{\setCJKsansfont{SimHei}}
    {\setCJKsansfont{FandolHei-Regular.otf}}}
\IfFontExistsTF{Times New Roman}{\setmainfont{Times New Roman}}
  {\setmainfont{Latin Modern Roman}}

\begin{document}

\begin{tikzpicture}[
  font=\small,
  node distance=0pt,
  % 阶段标题栏
  hbar/.style={draw,rounded corners=1.2pt,align=center,inner xsep=3pt,inner ysep=3.0pt,
               minimum width=62mm,fill=black!14,line width=0.7pt,
               font=\sffamily\small\bfseries},
  % 子步骤
  pbox/.style={draw,rounded corners=1.2pt,align=center,inner xsep=2pt,inner ysep=3.0pt,
               text width=58mm,fill=white,line width=0.5pt},
  % 阶段输出
  obox/.style={draw,rounded corners=1.2pt,align=center,inner xsep=2pt,inner ysep=3.0pt,
               text width=58mm,fill=black!7,line width=0.7pt},
  % 起点/终点
  tbox/.style={draw,rounded corners=5pt,align=center,inner xsep=2pt,inner ysep=3.0pt,
               text width=62mm,fill=black!4,line width=0.7pt},
  ar/.style={-{Latex[length=1.7mm,width=1.4mm]},line width=0.55pt},
]

% ============ 输入 ============
\node[tbox] (in) {候选字段集合 $F$ 与受保护目标 $s$};

% ============ 阶段一（3.2 节）============
\node[hbar,below=4.0mm of in]  (hA) {阶段一\quad 显式推断关系识别与去除};
\node[pbox,below=2.6mm of hA]  (a1) {已知公式剔除};
\node[pbox,below=2.6mm of a1]  (a2) {近似关系迭代去除};
\node[obox,below=2.6mm of a2]  (a3) {候选字段池 $\tilde F$（无显式推断路径）};

% ============ 阶段二（3.3 节）============
\node[hbar,below=4.6mm of a3]  (hB) {阶段二\quad 隐性推断源定位};
\node[pbox,below=2.6mm of hB]  (b1) {分解式稀疏门控网络训练};
\node[pbox,below=2.6mm of b1]  (b2) {活跃阈值 $\tau_g$ 筛选与独立重训验证};
\node[obox,below=2.6mm of b2]  (b3) {隐性推断源 $A^{\ast}$ 与风险度量 $(|A^{\ast}|,\rho)$};

% ============ 阶段三（3.4 节）：一个总阶段 ============
\node[hbar,below=4.6mm of b3]  (hC) {阶段三\quad 处置与残余风险再评估};

% ============ 输出 ============
\node[tbox,below=4.0mm of hC]  (fin) {通过审查，按处置后口径发布};

% ============ 箭头 ============
\foreach \x/\y in {in/hA,hA/a1,a1/a2,a2/a3,a3/hB,hB/b1,b1/b2,b2/b3,b3/hC,hC/fin}
  \draw[ar] (\x)--(\y);

% ============ 阶段一、二的虚线分组框（含标题栏）============
\begin{scope}[on background layer]
  \node[fit=(hA)(a3),draw=black!60,dashed,line width=0.55pt,rounded corners=2pt,
        inner xsep=1.8mm,inner ysep=1.5mm,fill=black!3] (gA) {};
  \node[fit=(hB)(b3),draw=black!60,dashed,line width=0.55pt,rounded corners=2pt,
        inner xsep=1.8mm,inner ysep=1.5mm,fill=black!3] (gB) {};
\end{scope}

\end{tikzpicture}

\end{document}
